% !TEX program = lualatex   % Рекомендуется для Metropolis; можно заменить на pdflatex
% Компиляция на Overleaf: Menu → Compiler → LuaLaTeX (или pdfLaTeX)
% Тема metropolis доступна на Overleaf из коробки.
% Альтернатива без внешних пакетов: закомментируйте metropolis и раскомментируйте Madrid.

\documentclass[aspectratio=169,10pt,russian]{beamer}

% --- Тема оформления ---
\usetheme{metropolis}
\metroset{progressbar=frametitle, subsectionpage=progressbar, numbering=fraction}
\usecolortheme{default}

% Запасной вариант (встроенная тема Beamer, без доп. пакетов):
% \usetheme{Madrid}
% \usecolortheme{whale}

% --- Язык и кодировка ---
\usepackage[T2A]{fontenc}
\usepackage[utf8]{inputenc}
\usepackage[russian]{babel}

% --- Полезные пакеты ---
\usepackage{booktabs}
\usepackage{tabularx}
\usepackage{multirow}
\usepackage{hyperref}
\hypersetup{
  colorlinks=true,
  linkcolor=structure.fg,
  urlcolor=structure.fg,
}

% --- Метаданные ---
\title{ИИ-инструменты: обзор, сравнение и практика}
\subtitle{От ChatGPT до агентов и локальных моделей}
\author{Имя Фамилия}          % <-- исправлено: убран комментарий внутри аргумента
\institute{Организация / курс}% <-- исправлено
\date{\today}

% Маркер «заполнить» для скелета
\newcommand{\todo}[1]{\textcolor{red!70!black}{\textit{[#1]}}}

\begin{document}

% ============================================================
\maketitle

\begin{frame}{О презентации}
  \begin{itemize}
    \item Цель: систематизировать опыт работы с ИИ-инструментами с конца 2022~г.
    \item Охват: типы решений, бенчмарки, доступность из РФ, сценарии и техники.
    \item Формат: обзор + личные выводы по \todo{реально протестированным} инструментам.
  \end{itemize}
\end{frame}

\begin{frame}{Содержание}
  \tableofcontents
\end{frame}

% ============================================================
\section{Временная шкала развития ИИ-инструментов}

\begin{frame}{Ключевые вехи (конец 2022 — настоящее время)}
  \small
  \begin{tabularx}{\linewidth}{@{}l X@{}}   % исправлено: добавлена закрывающая }
    \toprule
    \textbf{Период} & \textbf{Событие} \\
    \midrule
    Ноя 2022       & ChatGPT — массовый доступ к диалоговому ИИ \\
    Мар 2023       & GPT-4; начало «гонки» мультимодальных моделей \\
    Июл 2023       & Claude~2; расширение контекста и безопасности \\
    Ноя 2023       & GPT-4 Turbo, Assistants API; Gemini 1.0 \\
    2024           & GPT-4o, Claude~3, Llama~3; IDE-ассистенты (Cursor, Copilot) \\
    2024–2025      & Reasoning-модели (o1, o3); взрыв агентских фреймворков \\
    2025           & \todo{DeepSeek, Claude~4, GPT-4.1/5 — актуализировать} \\
    \bottomrule
  \end{tabularx}

  \vspace{0.5em}
  \todo{Добавить 2–3 лично значимых события (релиз протестированного инструмента)}
\end{frame}

\begin{frame}{Эволюция парадигмы использования}
  \begin{columns}[T,onlytextwidth]
    \column{0.48\textwidth}
    \textbf{2022–2023}
    \begin{itemize}
      \item Чат как «умный поиск»
      \item Копипаст кода в браузер
      \item Первые API-интеграции
    \end{itemize}

    \column{0.48\textwidth}
    \textbf{2024–2026}
    \begin{itemize}
      \item IDE и агенты «внутри» рабочего процесса
      \item Локальные и гибридные схемы
      \item Мультиагентные пайплайны
    \end{itemize}
  \end{columns}

  \vspace{1em}
  \centering
  \todo{Вставить простую timeline-диаграмму (TikZ или рисунок)}
\end{frame}

% ============================================================
\section{Типы ИИ-инструментов}

\begin{frame}{Карта экосистемы}
  \begin{itemize}
    \item \textbf{Чат-интерфейсы} — ChatGPT, Claude, Gemini, Perplexity\ldots
    \item \textbf{API} — OpenAI, Anthropic, Google, OpenRouter, Together\ldots
    \item \textbf{IDE / редакторы} — Cursor, GitHub Copilot, JetBrains AI, Continue
    \item \textbf{Агенты} — Devin-подобные, AutoGPT, CrewAI, Cursor Agent, Claude Code
    \item \textbf{Локальные модели} — Ollama, LM Studio, llama.cpp, vLLM
    \item \textbf{Российские платформы} — GigaChat, YandexGPT, Kandinsky, \todo{другие}
  \end{itemize}
\end{frame}

\begin{frame}{Чат vs API vs IDE vs агенты}
  \footnotesize
  \begin{tabularx}{\linewidth}{@{}l X X X@{}}   % исправлено
    \toprule
    \textbf{Тип} & \textbf{Сильные стороны} & \textbf{Слабые стороны} & \textbf{Когда выбирать} \\
    \midrule
    Чат   & Низкий порог входа & Мало контекста проекта & Быстрые вопросы, черновики \\
    API   & Автоматизация, масштаб & Нужна разработка & Скрипты, боты, пайплайны \\
    IDE   & Контекст кода, diff & Привязка к редактору & Ежедневная разработка \\
    Агенты & Многошаговые задачи & Непредсказуемость, стоимость & Исследования, прототипы \\
    Локальные & Приватность, офлайн & Железо, качество & Чувствительные данные \\
    \bottomrule
  \end{tabularx}
\end{frame}

\begin{frame}{Российские платформы — краткий обзор}
  \begin{itemize}
    \item \textbf{GigaChat (Сбер)} — \todo{модели, API, лимиты, ваш опыт}
    \item \textbf{YandexGPT / YandexART} — \todo{интеграции, качество для русского}
    \item \textbf{Kandinsky} — генерация изображений
    \item \textbf{Другие} — \todo{Visper, RuGPT, корпоративные решения}
  \end{itemize}

  \vspace{0.5em}
  \todo{Таблица: платформа / доступ без VPN / бесплатный tier / API}
\end{frame}

% ============================================================
\section{Сравнение протестированных решений}

\begin{frame}{Методология сравнения}
  \begin{itemize}
    \item \textbf{SWE-bench} — бенчмарк решения реальных GitHub-issues (код)
    \item \textbf{Chatbot Arena (LMSYS)} — слепое сравнение моделей пользователями (Elo)
    \item \textbf{Собственные тесты} — \todo{опишите 3–5 типовых задач}
  \end{itemize}

  \vspace{0.5em}
  \small
  Источники: \href{https://www.swebench.com/}{swebench.com},
  \href{https://chat.lmsys.org/}{Chatbot Arena}
\end{frame}

\begin{frame}{SWE-bench: результаты (шаблон)}
  \footnotesize
  \begin{tabular}{lrrl}
    \toprule
    \textbf{Инструмент / модель} & \textbf{Resolved \%} & \textbf{Verified} & \textbf{Комментарий} \\
    \midrule
    \todo{Model A}  & -- & \todo{да/нет} & \todo{} \\
    \todo{Model B}  & -- & \todo{да/нет} & \todo{} \\
    \todo{Model C}  & -- & \todo{да/нет} & \todo{} \\
    \todo{Cursor Agent} & -- & -- & \todo{личный опыт} \\
    \bottomrule
  \end{tabular}

  \vspace{0.5em}
  \todo{Подставить актуальные цифры на дату презентации}
\end{frame}

\begin{frame}{Chatbot Arena: лидеры (шаблон)}
  \footnotesize
  \begin{tabular}{lrl}
    \toprule
    \textbf{Модель} & \textbf{Arena Elo} & \textbf{Заметки} \\
    \midrule
    \todo{Top-1} & ---- & \todo{} \\
    \todo{Top-2} & ---- & \todo{} \\
    \todo{Top-3} & ---- & \todo{} \\
    \todo{Протестированная вами} & ---- & \todo{субъективная оценка} \\
    \bottomrule
  \end{tabular}

  \vspace{0.5em}
  \todo{Скриншот или ссылка на актуальный leaderboard}
\end{frame}

\begin{frame}{Сводная таблица «что выбрал бы для…»}
  \footnotesize
  \begin{tabularx}{\linewidth}{@{}l XXX@{}}   % исправлено
    \toprule
    \textbf{Задача} & \textbf{1-й выбор} & \textbf{Запасной} & \textbf{Не рекомендую} \\
    \midrule
    Написание кода     & \todo{} & \todo{} & \todo{} \\
    Рефакторинг        & \todo{} & \todo{} & \todo{} \\
    Тексты / перевод   & \todo{} & \todo{} & \todo{} \\
    Исследование темы  & \todo{} & \todo{} & \todo{} \\
    Работа с документами & \todo{} & \todo{} & \todo{} \\
    \bottomrule
  \end{tabularx}
\end{frame}

% ============================================================
\section{Доступность из России}

\begin{frame}{Ограничения и обходные пути}
  \begin{itemize}
    \item \textbf{Геоблокировки} — ChatGPT, Claude, часть API; статус на \todo{дату}
    \item \textbf{Оплата} — зарубежные карты, посредники, корпоративные аккаунты
    \item \textbf{VPN / прокси} — риски ToS, стабильность, latency
    \item \textbf{Локальные альтернативы} — без VPN, но другой уровень качества
    \item \textbf{OpenRouter / прокси-API} — \todo{ваш опыт}
  \end{itemize}
\end{frame}

\begin{frame}{Матрица доступности (шаблон)}
  \footnotesize
  \begin{tabularx}{\linewidth}{@{}l X X X@{}}   % исправлено
    \toprule
    \textbf{Сервис} & \textbf{Без VPN} & \textbf{Оплата из РФ} & \textbf{API} \\
    \midrule
    ChatGPT / OpenAI & \todo{} & \todo{} & \todo{} \\
    Claude           & \todo{} & \todo{} & \todo{} \\
    Gemini           & \todo{} & \todo{} & \todo{} \\
    Cursor           & \todo{} & \todo{} & \todo{} \\
    GigaChat         & \todo{} & \todo{} & \todo{} \\
    YandexGPT        & \todo{} & \todo{} & \todo{} \\
    Ollama (локально)& Да & N/A & Локально \\
    \bottomrule
  \end{tabularx}

  \vspace{0.5em}
  \small
  \todo{Пометить ✓/✗/частично; добавить юридический дисклеймер при необходимости}
\end{frame}

% ============================================================
\section{Примеры использования}

\begin{frame}{В работе}
  \begin{itemize}
    \item Code review и генерация тестов — \todo{пример}
    \item Документация и README — \todo{пример}
    \item Отладка и объяснение stack trace — \todo{пример}
    \item Автоматизация рутины (скрипты, CI) — \todo{пример}
  \end{itemize}
\end{frame}

\begin{frame}{В учёбе}
  \begin{itemize}
    \item Разбор сложных тем «на пальцах» — \todo{предмет}
    \item Подготовка к экзаменам, карточки — \todo{}
    \item Проверка решений (не подмена обучения!) — \todo{этика}
    \item Работа с литературой: summary, связи — \todo{}
  \end{itemize}
\end{frame}

\begin{frame}{Планирование и исследования}
  \begin{itemize}
    \item \textbf{Планирование} — декомпозиция проектов, оценка рисков, чек-листы
    \item \textbf{Исследования} — Perplexity / Deep Research / агентный поиск
    \item \textbf{Ограничения} — галлюцинации, устаревшие данные, нужна верификация
    \item \todo{1 кейс: «как решал(а) задачу от постановки до результата»}
  \end{itemize}
\end{frame}

% ============================================================
\section{Техники работы с ИИ}

\begin{frame}{Промптинг и структура запросов}
  \begin{itemize}
    \item \textbf{Rôle + Context + Task + Format} — базовый шаблон
    \item \textbf{Few-shot} — примеры желаемого вывода
    \item \textbf{Chain-of-Thought} — «рассуждай по шагам» для сложной логики
    \item \textbf{Итеративное уточнение} — черновик → критика → доработка
  \end{itemize}
\end{frame}

\begin{frame}{Контекст, RAG и инструменты}
  \begin{itemize}
    \item \textbf{Контекстное окно} — что включать: файлы, diff, логи
    \item \textbf{RAG} — база знаний, embeddings, chunking
    \item \textbf{Tool use / function calling} — калькулятор, поиск, API
    \item \textbf{Агентные паттерны} — planner–executor, рефлексия, human-in-the-loop
  \end{itemize}

  \vspace{0.5em}
  \todo{Мини-демо: один «хороший» и один «плохой» промпт}
\end{frame}

\begin{frame}{Антипаттерны}
  \begin{itemize}
    \item Слепое доверие без проверки фактов и кода
    \item Слишком длинный промпт без структуры
    \item Передача секретов и PII в облачные сервисы
    \item Ожидание «магии» от агента без постановки критериев успеха
  \end{itemize}
\end{frame}

% ============================================================
\section{Практические выводы}

\begin{frame}{Что оказалось полезным}
  \begin{itemize}
    \item \todo{Инструмент 1 — конкретная польза}
    \item \todo{Инструмент 2 — конкретная польза}
    \item \todo{Приём / workflow, который стабильно работает}
    \item \todo{Комбинация инструментов (например, IDE + чат + локальная модель)}
  \end{itemize}
\end{frame}

\begin{frame}{Что неудобно или переоценено}
  \begin{itemize}
    \item \todo{Высокая стоимость / нестабильный доступ}
    \item \todo{Агенты: где тратят время впустую}
    \item \todo{Модели, которые хуже ожиданий на ваших задачах}
    \item \todo{Скрытые costs: время на промпты, review, исправление ошибок}
  \end{itemize}
\end{frame}

\begin{frame}{Рекомендации аудитории}
  \begin{enumerate}
    \item Начните с \todo{одного чата + одного IDE-ассистента}.
    \item Заведите личный набор «эталонных задач» для сравнения моделей.
    \item Для чувствительных данных — \todo{локальная или российская схема}.
    \item Всегда верифицируйте код, факты и ссылки.
    \item \todo{Ваш главный вывод в одном предложении}
  \end{enumerate}
\end{frame}

% ============================================================
\begin{frame}[standout]
  Вопросы?

  \vspace{1em}
  \small
  \todo{Контакты, ссылки на материалы, репозиторий с промптами}
\end{frame}

\appendix

\begin{frame}{Приложение: полезные ссылки}
  \small
  \begin{itemize}
    \item SWE-bench — \url{https://www.swebench.com/}
    \item Chatbot Arena — \url{https://chat.lmsys.org/}
    \item Papers With Code — \url{https://paperswithcode.com/}
    \item \todo{Ваши избранные блоги, документация, каналы}
  \end{itemize}
\end{frame}

\begin{frame}{Приложение: глоссарий}
  \footnotesize
  \begin{tabularx}{\linewidth}{@{}l X@{}}   % исправлено
    \toprule
    \textbf{Термин} & \textbf{Кратко} \\
    \midrule
    LLM & Large Language Model — большая языковая модель \\
    RAG & Retrieval-Augmented Generation \\
    Elo & Рейтинговая система Chatbot Arena \\
    SWE-bench & Бенчмарк программирования на реальных issues \\
    Agent & Система с планированием и вызовом инструментов \\
    \bottomrule
  \end{tabularx}
\end{frame}

\end{document}